\section{Regime-local Mixture-of-Experts reduced-order model specialists}
\label{sec:hprs_specialists}





Each regime-local ROM chart $\mathcal{C}^{\mathrm{ROM}}_r$ is augmented with a
regime-specific reduced dynamical model, referred to below as a
\emph{regime-local specialist}.  The three specialists share the same hybrid
physics--data and sparse mixture-of-experts architecture, 
but differ in the underlying flow physics.  Within a specialist, the term
\emph{expert} denotes one neural subnetwork of the internal sparse mixture;
it does not refer to the complete regime-local ROM.  This terminology
distinguishes the expert-level sparsity inside each specialist from the
inter-regime routing and physical-field combination introduced in
Section~\ref{sec:hierarchical_ral_moe}.

\subsection{Semi-explicit hybrid physics--data reduced dynamics}
\label{subsec:semi_explicit_hybrid_dynamics}

For a fixed regime $r\in\mathcal{R}$, the velocity coefficients are governed
by a Galerkin reduced operator supplemented by a learned closure term,
\begin{equation}
  \dot{\boldsymbol{a}}^{(r)}
  =
  \mathscr{F}^{\mathrm{G}}_r
  \bigl(
    \boldsymbol{a}^{(r)},
    \boldsymbol{b}^{(r)};
    \mu
  \bigr)
  +
  \boldsymbol{\rho}^{u}_{\theta_r}
  \bigl(
    \boldsymbol{\xi}^{(r)}
  \bigr).
  \label{eq:hybrid_velocity_rhs}
\end{equation}
Here $\mathscr{F}^{\mathrm{G}}_r$ is the regime-local POD--Galerkin operator
defined in \eqref{eq:chart_local_galerkin_rhs_definition}, and
$\boldsymbol{\rho}^{u}_{\theta_r}$ approximates the component of the reduced
velocity dynamics that is not represented adequately by the projected
operator.  In the notation of the preceding subsection,
\begin{equation}
  \mathscr{F}^{\mathrm{G}}_r
  (\boldsymbol{a},\boldsymbol{b};\mu)
  =
  \boldsymbol{c}_{r}(\mu)
  +A_{r}(\mu)\boldsymbol{a}
  +H_{r}(\mu):
   \bigl(\boldsymbol{a}\otimes\boldsymbol{a}\bigr)
  +C_{r}^{p}(\mu)\boldsymbol{b}.
  \label{eq:galerkin_rhs}
\end{equation}
Thus, the learned term acts as a closure correction to an explicitly retained
physics-based reduced operator rather than replacing the Galerkin dynamics.

Let $\Delta t_r(\mu)$ denote the sampling interval associated with the
regime-$r$ trajectory at parameter value $\mu$.  A complete velocity
macro-step is written abstractly as
\begin{equation}
  \boldsymbol{a}^{(r)}_{n+1}
  =
  \mathcal{I}_{r}
  \left(
    \boldsymbol{a}^{(r)}_{n};
    \mathscr{F}^{\mathrm{G}}_r
    +\boldsymbol{\rho}^{u}_{\theta_r},
    \Delta t_r(\mu),
    \mathcal{C}^{(r)}_n
  \right),
  \label{eq:native_velocity_integrator}
\end{equation}
where $\mathcal{I}_{r}$ is the time-integration operator used by the
specialist and $\mathcal{C}^{(r)}_n$ denotes the auxiliary context held fixed
or prescribed during the macro-step, including the pressure coefficients,
history variables, parameter value, and any external temporal descriptor.
The time increment is an argument of the numerical integrator and is not
included as a trainable network feature.  All three specialists employ the classical fourth-order Runge–Kutta method.

Pressure is not assigned an independent learned evolution equation.
Instead, after the velocity update, the regime-local Pressure--Poisson map
defined in \eqref{eq:chart_local_pressure_poisson_map} provides a
physics-based estimate
\begin{equation}
  \boldsymbol{b}^{PP,(r)}_{n+1}
  =
  \mathscr{Q}^{PP}_{r}
  \bigl(
    \boldsymbol{a}^{(r)}_{n+1};
    \mu
  \bigr).
  \label{eq:pressure_poisson_base}
\end{equation}
To complement this prediction, each specialist $r\in\{S,H,P\}$ also
supplies a data-driven candidate and a componentwise gate through the
learned networks $\boldsymbol{\rho}^{p}_{\theta_r}$ and
$\boldsymbol{\eta}^{p}_{\theta_r}$:
\begin{align}
  \widehat{\boldsymbol{b}}^{D,(r)}_{n+1}
  &=
  \boldsymbol{\rho}^{p}_{\theta_r}
  \bigl(
    \boldsymbol{\xi}^{(r)}_n
  \bigr),
  \label{eq:data_driven_pressure_candidate}\\
  \boldsymbol{\alpha}^{(r)}_n
  &=
  \operatorname{sigmoid}
  \left(
    \boldsymbol{\eta}^{p}_{\theta_r}
    \bigl(
      \boldsymbol{\xi}^{(r)}_n
    \bigr)
  \right),
  \qquad
  \boldsymbol{\alpha}^{(r)}_n\in(0,1)^{d_p}.
  \label{eq:adaptive_pressure_gate}
\end{align}
The final pressure coefficients are then obtained as a componentwise
convex combination of the physics-based estimate and the data-driven
candidate:
\begin{equation}
  \boldsymbol{b}^{(r)}_{n+1}
  =
  \boldsymbol{\alpha}^{(r)}_n
  \odot
  \boldsymbol{b}^{PP,(r)}_{n+1}
  +
  \bigl(
    \boldsymbol{1}-\boldsymbol{\alpha}^{(r)}_n
  \bigr)
  \odot
  \widehat{\boldsymbol{b}}^{D,(r)}_{n+1},
  \qquad
  r\in\{S,H,P\}.
  \label{eq:adaptive_pressure_closure}
\end{equation}


Equations \eqref{eq:hybrid_velocity_rhs}--\eqref{eq:adaptive_pressure_closure}
therefore define a partitioned semi-explicit reduced model: the velocity
coefficients are advanced by a differential update, whereas the pressure
coefficients are reconstructed algebraically after each velocity macro-step.

\subsection{Feature representation and internal sparse mixture-of-experts architecture}
\label{subsec:sparse_moe_features}
\subsection{Feature representation and internal sparse mixture-of-experts architecture}
\label{subsec:sparse_moe_features}

The learned closure is evaluated from a regime-dependent feature vector that
combines the current reduced state, the Galerkin contribution, and the
physical parameter.  The current-state block is
\begin{equation}
  \boldsymbol{\chi}^{(r)}_n
  =
  \operatorname{col}
  \left(
    \widetilde{\boldsymbol{a}}^{(r)}_n,
    \widetilde{\boldsymbol{b}}^{(r)}_n,
    \widetilde{\boldsymbol{g}}^{(r)}_n,
    \widetilde{\mu}
  \right),
  \label{eq:current_feature_block}
\end{equation}
where
\begin{equation}
  \boldsymbol{g}^{(r)}_n
  =
  \mathscr{F}^{\mathrm{G}}_r
  \bigl(
    \boldsymbol{a}^{(r)}_n,
    \boldsymbol{b}^{(r)}_n;
    \mu
  \bigr).
  \label{eq:current_galerkin_feature}
\end{equation}
The quantities $\widetilde{\boldsymbol{a}}^{(r)}_n$ and
$\widetilde{\boldsymbol{b}}^{(r)}_n$ are standardized using the
regime-local statistics defined in \eqref{eq:local_scalers}.  The Galerkin
feature $\widetilde{\boldsymbol{g}}^{(r)}_n$ and the scalar parameter
feature $\widetilde{\mu}$ are standardized analogously using statistics
computed from the same regime-local training population.

Each learned closure uses the current state together with the two preceding
states.  For $\ell\in\{1,2\}$, the history block is constructed analogously
from the corresponding time step,
\begin{equation}
  \boldsymbol{h}^{(r)}_{n,\ell}
  =
  \operatorname{col}
  \left(
    \widetilde{\boldsymbol{a}}^{(r)}_{n-\ell},
    \widetilde{\boldsymbol{b}}^{(r)}_{n-\ell},
    \widetilde{\boldsymbol{g}}^{(r)}_{n-\ell},
    \widetilde{\mu}
  \right).
  \label{eq:history_feature_block}
\end{equation}
The complete common feature vector is
\begin{equation}
  \boldsymbol{\xi}^{(r)}_n
  =
  \operatorname{col}
  \left(
    \boldsymbol{\chi}^{(r)}_n,
    \boldsymbol{h}^{(r)}_{n,1},
    \boldsymbol{h}^{(r)}_{n,2}
  \right).
  \label{eq:complete_specialist_features}
\end{equation}


The learned closure uses hierarchical sparse routing inside each specialist.
Let $\mathcal{G}$ denote the set of expert groups and let
$\boldsymbol{\pi}^{\mathrm{grp}}(\boldsymbol{\xi})$ be the normalized output
of the group router.  During sparse evaluation, one group is selected by
\begin{equation}
  g^{\star}
  =
  \underset{g\in\mathcal{G}}{\operatorname{arg\,max}}
  \;\pi^{\mathrm{grp}}_{g}
  \bigl(
    \boldsymbol{\xi}
  \bigr).
  \label{eq:group_top1}
\end{equation}
Within the selected group, velocity and pressure use separate channel
routers.  For channel $c\in\{u,p\}$, let
$\mathcal{T}^{c}_{2}(\boldsymbol{\xi};g^{\star})$ contain the two routed
experts with the largest channel-specific routing scores.  Each group also
contains one shared expert, indexed by $e=0$, that is evaluated for every
input.  The channel output is
\begin{equation}
  \boldsymbol{\mathcal{M}}^{c}_{\theta_r}
  \bigl(
    \boldsymbol{\xi}
  \bigr)
  =
  \omega^{c}_{0}
  \mathcal{E}^{c}_{0,g^{\star}}
  \bigl(
    \boldsymbol{\xi}
  \bigr)
  +
  \sum_{e\in\mathcal{T}^{c}_{2}}
  \omega^{c}_{e}
  \mathcal{E}^{c}_{e,g^{\star}}
  \bigl(
    \boldsymbol{\xi}
  \bigr),
  \label{eq:channel_sparse_moe}
\end{equation}
where
\begin{equation}
  \omega^{c}_{e}\geq 0,
  \qquad
  \omega^{c}_{0}
  +
  \sum_{e\in\mathcal{T}^{c}_{2}}
  \omega^{c}_{e}
  =
  1.
  \label{eq:channel_sparse_weights}
\end{equation}
Thus, one group, one shared expert, and two routed experts are active for each
output channel.  Separate velocity and pressure routers are retained because
the two channels approximate different closure quantities: the velocity
channel corrects the reduced differential operator, whereas the pressure
channel supplies either an algebraic correction or a data-driven pressure
candidate.  Accordingly,
$\boldsymbol{\mathcal{M}}^{u}_{\theta_r}$ represents
$\boldsymbol{\rho}^{u}_{\theta_r}$, while the interpretation of
$\boldsymbol{\mathcal{M}}^{p}_{\theta_r}$ follows
\eqref{eq:steady_pressure_closure} or
\eqref{eq:data_driven_pressure_candidate}, depending on the regime.

Each expert combines a nonlinear multilayer branch, an affine branch, and a
low-rank bilinear branch.  For an input dimension $m_r$ and an output
dimension $d_c$, the expert map is
\begin{equation}
  \mathcal{E}^{c}_{e,g}
  (\boldsymbol{\xi})
  =
  \mathcal{N}^{c}_{e,g}
  (\boldsymbol{\xi})
  +
  W^{c,\mathrm{lin}}_{e,g}
  \boldsymbol{\xi}
  +
  \kappa
  W^{c,o}_{e,g}
  \left[
    \bigl(
      {U^{c}_{e,g}}^{\mathsf T}\boldsymbol{\xi}
    \bigr)
    \odot
    \bigl(
      {V^{c}_{e,g}}^{\mathsf T}\boldsymbol{\xi}
    \bigr)
  \right],
  \label{eq:physics_aware_expert}
\end{equation}
with
\begin{equation}
  U^{c}_{e,g},V^{c}_{e,g}
  \in
  \mathbb{R}^{m_r\times q_b},
  \qquad
  W^{c,o}_{e,g}
  \in
  \mathbb{R}^{d_c\times q_b},
  \qquad
  q_b=4,
  \qquad
  \kappa=0.05.
  \label{eq:expert_bilinear_dimensions}
\end{equation}


The bilinear branch represents a rank-$q_b$
approximation of a quadratic form of the input, providing a low-dimensional mechanism for capturing pairwise interactions among reduced coordinates and physical features. Its presence is motivated by the polynomial structure of the projected convective term, but it does not impose the Galerkin tensor itself. The bilinear branch uses a rank of $q_b=4$, which provides sufficient
capacity to represent the dominant pairwise interactions among the
$m_r$ input features while keeping the parameter count an order of
magnitude below that of a full quadratic form.
The scalar $\kappa=0.05$ is a small pre-factor that stabilizes training
by damping the contribution of the bilinear branch in the initial
optimization phase and by preventing over-reliance on quadratic
corrections to the already physics-informed operator.


\subsection{Regime-specific dynamical formulations}
\label{subsec:regime_specific_contracts}

Although the three specialists share the same hybrid architecture and
feature representation, each is trained on a distinct regime‑local
set and deploys its own parameter set $\theta_r$.  All specialists
employ the classical fourth‑order Runge–Kutta (RK4) method for the
velocity update and the gated pressure closure
\eqref{eq:adaptive_pressure_closure}.  The feature vector
$\boldsymbol{\xi}^{(r)}_n$ is constructed from the current and history
blocks as described in~\S\ref{subsec:sparse_moe_features}.

\paragraph{Velocity update.}
During one macro-step, the pressure coefficients $\boldsymbol{b}^{(r)}_n$,
the history blocks, and the physical parameter $\mu$ are held frozen.
Let $\mathcal{C}^{(r)}_n$ denote this frozen context.  The right‑hand side
of the velocity equation is
\begin{equation}
  F_{r}
  \bigl(
    \boldsymbol{a};
    \mathcal{C}^{(r)}_n
  \bigr)
  =
  \mathscr{F}^{\mathrm{G}}_{r}
  \bigl(
    \boldsymbol{a},
    \boldsymbol{b}^{(r)}_n;
    \mu
  \bigr)
  +
  \boldsymbol{\rho}^{u}_{\theta_r}
  \bigl(
    \boldsymbol{\xi}^{(r)}(\boldsymbol{a};\mathcal{C}^{(r)}_n)
  \bigr),
  \label{eq:frozen_context_rhs}
\end{equation}
where $\boldsymbol{\xi}^{(r)}(\boldsymbol{a};\mathcal{C}^{(r)}_n)$
denotes the feature vector evaluated with the velocity coefficient
$\boldsymbol{a}$ and the frozen context.  The four RK4 stages are
\begin{align}
  \boldsymbol{k}_1
  &= F_{r}\bigl( \boldsymbol{a}^{(r)}_n; \mathcal{C}^{(r)}_n \bigr),
  \nonumber\\
  \boldsymbol{k}_2
  &= F_{r}\left( \boldsymbol{a}^{(r)}_n + \tfrac{\Delta t_r(\mu)}{2}\boldsymbol{k}_1; \mathcal{C}^{(r)}_n \right),
  \nonumber\\
  \boldsymbol{k}_3
  &= F_{r}\left( \boldsymbol{a}^{(r)}_n + \tfrac{\Delta t_r(\mu)}{2}\boldsymbol{k}_2; \mathcal{C}^{(r)}_n \right),
  \nonumber\\
  \boldsymbol{k}_4
  &= F_{r}\bigl( \boldsymbol{a}^{(r)}_n + \Delta t_r(\mu)\boldsymbol{k}_3; \mathcal{C}^{(r)}_n \bigr),
  \nonumber\\
  \boldsymbol{a}^{(r)}_{n+1}
  &=
  \boldsymbol{a}^{(r)}_n
  +
  \frac{\Delta t_r(\mu)}{6}
  \bigl(
    \boldsymbol{k}_1 + 2\boldsymbol{k}_2 + 2\boldsymbol{k}_3 + \boldsymbol{k}_4
  \bigr).
  \label{eq:unified_rk4}
\end{align}
After the velocity update, the pressure coefficients are obtained through
the gated combination
\eqref{eq:adaptive_pressure_closure}, using the new velocity
coefficients $\boldsymbol{a}^{(r)}_{n+1}$ to compute the physics‑based
estimate $\boldsymbol{b}^{PP,(r)}_{n+1}$ and evaluating the data‑driven
candidate and gate networks on the feature vector
$\boldsymbol{\xi}^{(r)}_n$ (which still contains the pre‑update
pressure).  The history blocks are then advanced by one step, and the
procedure repeats.

\paragraph{Training.}
All three specialists are trained with a progressive closed‑loop rollout
curriculum, identical in form to the one originally used for the
periodic regime.  At each curriculum stage, the model predicts a
trajectory of length $K\in\{4,8,12,16\}$ steps from an initial condition
drawn from the regime‑specific training set $\mathfrak{D}_{r}^{\mathrm{tr}}$.
The loss function is the mean squared error between the predicted and
the reference POD coefficients, normalized by the regime‑local standard
deviations:
\begin{equation}
  \mathcal{L}_{r}
  =
  \frac{1}{K}\sum_{k=1}^{K}
  \left[
    \frac{1}{d_u}
    \bigl\|
      \bigl(
        \boldsymbol{a}^{(r)}_{n+k} - \widehat{\boldsymbol{a}}^{(r)}_{n+k}
      \bigr)
      \oslash
      \boldsymbol{\sigma}^{(r)}_a
    \bigr\|_2^2
    +
    \frac{1}{d_p}
    \bigl\|
      \bigl(
        \boldsymbol{b}^{(r)}_{n+k} - \widehat{\boldsymbol{b}}^{(r)}_{n+k}
      \bigr)
      \oslash
      \boldsymbol{\sigma}^{(r)}_b
    \bigr\|_2^2
  \right].
  \label{eq:unified_rollout_loss}
\end{equation}
Here hatted quantities denote reference trajectories from the training
data, and $\boldsymbol{\sigma}^{(r)}_a$, $\boldsymbol{\sigma}^{(r)}_b$ are
the regime‑local coefficient standard deviations defined in
\eqref{eq:local_scalers}.  No additional contraction, fixed‑point, or
radial regularization terms are introduced; regime specialization relies
entirely on the distinct training populations and the resulting
parameter sets $\theta_r$.





Model selection is performed using complete parameter trajectories
assigned to the validation set.  After validation‑based model selection, the parameters and all
regime‑specific preprocessing and reduced operators are frozen.
Subsequent optimization of the inter‑regime router or any admissible
overlap‑fusion model does not update the local POD means and bases,
normalization statistics, Galerkin and Pressure–Poisson operators,
internal routers, neural experts, time‑integration schemes, or
pressure‑reconstruction maps.  The resulting hierarchical ROM is
therefore assembled from independently trained and validation‑selected
local dynamical models. 
